% Original pages: 337--341
% Figures: none.
% Uncertainty log: final sentence on original p. 341 truncates at the page boundary.

With this accomplished, evaluating $\operatorname{Tr}(-1)^F$ is as easy as in the
Wess-Zumino model. Only the ``vacuum states'' have $E = 0$. Other
states have energy at least equal to the lowest allowed momentum for
modes of non-zero momentum in the box of length $L$.

One can expand around $A_\mu^a = 0$; this ``vacuum'' contributes one to
$\operatorname{Tr}(-1)^F$. There are other vacuum states that can be obtained from
$A_\mu^a = 0$ by topologically non-trivial gauge transformations. In a
world defined by a given value of the vacuum angle $\theta$ there are still
$N$ sectors of configurations of $F_{\mu\nu}^a = 0$ (they are created from $A_\mu^a = 0$
by the gauge transformations which, according to 't Hooft, measure
the electric flux in the $z$ direction). So altogether, for $SU(N)$,
$\operatorname{Tr}(-1)^F = N$, and supersymmetry is not spontaneously broken.

Twisted boundary conditions may be introduced for any group
which has a non-trivial center. However, for groups other than $SU(N)$,
the twisted boundary conditions are not very useful in calculating
$\operatorname{Tr}(-1)^F$. The problem is that, for other groups, the twisted boundary
conditions do not eliminate the zero momentum mode because the
condition analogous to (41) does not imply $\phi_0 = 0$.

However, it is possible to calculate $\operatorname{Tr}(-1)^F$ with untwisted
boundary conditions, by quantizing the zero momentum modes in a
Born-Oppenheimer approximation. One finds\textsuperscript{12} that for a simple non-Abelian
Lie group of rank $r$, $\operatorname{Tr}(-1)^F = r + 1$. Thus, spontaneous supersymmetry
breaking does not occur in these theories.

What happens if additional matter fields are added? If the additional
fields are in a real representation of the gauge group, bare
masses $m_i$ are possible for all of the matter fields. In this case
all states containing the new quanta have non-zero energy, equal to
or greater than the smallest of the $m_i$. The new fields do not contribute
to $\operatorname{Tr}(-1)^F$. So supersymmetry is unbroken, just as if the
new fields were not present.

This argument shows that the ground state energy is zero for any
non-zero values of the $m_i$. Taking now the limit as $m_i \to 0$, the ground
state energy must remain zero. The only assumption needed here is that
the zero mass limit should exist. We conclude that supersymmetry is
unbroken for a theory with massless charged matter fields, as long as
they lie in a representation of the gauge group such that they could
have had bare masses.

In the very interesting case of theories with massless fields
in a complex representation of the gauge group --- so that gauge invariant
bare masses are impossible --- I do not know how to calculate
$\operatorname{Tr}(-1)^F$. For reasons explained elsewhere,\textsuperscript{13} there are difficulties
even in formulating the problem.

I should also mention that for a theory with gauge group $U(1)$,
$\operatorname{Tr}(-1)^F = 0$, as long as all charged fields have or could have had
bare masses. It is nonetheless possible to prove by a variant of
the concept of $\operatorname{Tr}(-1)^F$ that dynamical supersymmetry breaking does
not occur in any $U(1)$ gauge theory in which the Hamiltonian commutes
with charge conjugation invariance (or in which charge conjugation
invariance is broken only by Yukawa couplings, and not by the gauge
couplings or the Fayet-Iliopoulos $D$ term).

Likewise, in a theory with a simple gauge group that is spontaneously
broken at the tree level to a subgroup such as $SU(3) \times$
$SU(2) \times U(1)$ that contains a $U(1)$ factor, $\operatorname{Tr}(-1)^F = 0$.

Clearly, these results impose significant restrictions on the
possibility of supersymmetry breaking by strong gauge forces. However,
loopholes remain, such as the question of matter fields in a
complex representation of the gauge group. It also remains for the
future to determine whether these methods can shed light on the
apparent difficulties in $3+1$ dimensions in supersymmetry breaking by
instantons.

\section{ANOTHER APPROACH TO MASS HIERARCHIES}

In this last lecture I would like to describe a different approach
to the gauge hierarchy problem.

It is usually assumed, in thinking about the hierarchy problem,
that the ``large'' masses (the mass scale of grand unification and the
Planck mass) are the fundamental ones. The problem then is to explain
why the ``small'' masses (the masses of ordinary particle physics)
are so small. One may seek to explain this by finding a non-perturbative
mechanism that generates tiny masses. It was this possibility
that motivated the discussion of dynamical supersymmetry breaking in
the last few lectures.

Another possibility is to assume that the small masses are the
fundamental ones. One must then explain why the large masses are so
large.

This approach was actually followed by Dirac\textsuperscript{7} in his original
approach to the problem of the large numbers. Dirac assumed that
the electron and proton masses were the basic ones. To account for
the fact that the Planck mass is enormously larger, Dirac postulated
that the Planck mass is not constant in time but increases in time as
the universe grows older. The enormous present value of the Planck
mass was thus attributed by Dirac to the fact that (in elementary
particle units) the present universe is extremely old. Unfortunately,
this beautiful idea has some serious empirical difficulties,
which I will mention later.

Today I will be describing a class of supersymmetric theories\textsuperscript{14}
which spontaneously generate a mass scale vastly larger than the
mass scale postulated in the Lagrangian. In these theories it is
conceivable that the fundamental mass scale of nature could be comparable
to the mass scale of the Weinberg-Salam model. All larger
masses would be dynamically generated. As we will see, these
theories have something in common with Dirac's approach, and we can,
but need not, retrieve Dirac's idea that the ``constants'' of nature
are slowly varying functions of time.

In our previous discussions we emphasized dynamical (non-perturbative)
symmetry breaking. Now, however, we will consider theories in which
supersymmetry is spontaneously broken at the tree
level (classically). There will be no mystery about supersymmetry
breaking; we will simply choose a scalar potential whose minimum is
not invariant under supersymmetry. The problem will be to extract
the consequences of supersymmetry breaking.

To understand the idea, let us consider one of the simplest
models with supersymmetry spontaneously broken at the tree level ---
the O'Raifeartaigh model. As we discussed before, in this model
there are three complex fields $A$, $X$, and $Y$ of spin zero. The superspace
potential is $W(A,X,Y) = \lambda X(A^2 - M^2) + g\, Y\, A$, and the ordinary
potential energy is $V(A,X,Y) = |\partial W/\partial A|^2 + |\partial W/\partial X|^2 + |\partial W/\partial Y|^2$ or
\begin{equation}
V(A,X,Y) = \lambda^2 |A^2 - M^2|^2 + g^2 |A|^2 + |2 \lambda\, A X + gY|^2 .
\tag{42}
\end{equation}
This potential is strictly positive, so supersymmetry is spontaneously
broken.

We determine $A$ by minimizing the first two terms. For $(g/\lambda M)^2 > 2$
one finds $A = 0$; for $(g/\lambda M)^2 < 2$ one finds $A = \sqrt{M^2 - g^2/4\lambda^2}$. Supersymmetry
is spontaneously broken in either case.

Although minimization of the potential determines $A$, it does not
determine $X$ and $Y$ uniquely. Since the only term in the potential that
depends on $X$ and $Y$ is $|2 \lambda\, A X + g\, Y|^2$, we can choose any $X$, as long as
\begin{equation}
Y = -2\, \lambda\, A\, X / g
\tag{43}
\end{equation}
Classically, $X$ may be arbitrarily large; the energy is minimized as
long as (43) is satisfied.

This degeneracy is not a property of this one model alone.
Such degeneracies occur in many models in which supersymmetry is
spontaneously broken at the tree level.

Although the energy is independent of $X$, the particle masses
are not. A glance at (42) shows that for $X \gg M$ the mass of the $A$
particle is approximately $2 \lambda X$. If $X \gg M$, the theory has at least
two different energy scales. The small scale is the mass $M$ in the
Lagrangian. The large scale is the vacuum expectation value of $X$,
which is undetermined by the classical Lagrangian. In an extended
version of this model, we may try to interpret these as the scales
of weak interactions and of grand unification, respectively.

However, we should not \underline{arbitrarily} assume $X \gg M$. We should
\underline{determine} $X$ by calculating quantum corrections that remove the classical
degeneracy.

Let us recall how this goes.\textsuperscript{15} For bosons, the zero point
energy per unit volume is positive and has the form
\begin{equation}
\frac{1}{2}\,\hbar \sum_i \omega_i = \frac{1}{2}\,\hbar \sum_i \int \frac{d^3 k}{(2\pi)^3}\, \sqrt{k^2 + M_i^2(X)}
\tag{44}
\end{equation}
The sum on the right-hand side of (44) is a sum over the spin states
of the various bosons; the $M_i$ are the masses of the bosons, which
depend on $X$. For fermions we have instead the \underline{negative} energy from
filling the Dirac sea. It is
\begin{equation}
- \frac{1}{2}\,\hbar \sum_i \int \frac{d^3 k}{(2\pi)^3}\, \sqrt{k^2 + M_i^2(X)}
\tag{45}
\end{equation}
where now the sum runs over the fermion spin states.

If supersymmetry is unbroken at the tree level, the bosons and
fermions have the same masses. Then (44) and (45) cancel. This is
an illustration of the ``non-renormalization'' theorems: the ground
state energy is zero to all finite orders if it is zero classically.
We are, however, interested in cases in which supersymmetry is

% VERIFY: the final sentence is truncated at the original p. 341 page boundary and continues on p. 342, which is outside this assignment.
